\chapter*{Introducción}
\addcontentsline{toc}{chapter}{Introducción}

El auge del aprendizaje automático y en particular del aprendizaje profundo, ha transformado numerosas industrias y campos de investigación. Modelos como las redes neuronales artificiales han demostrado una capacidad sin precedentes para extraer patrones complejos de grandes volúmenes de datos, impulsando avances significativos en áreas como la visión por computadora, el procesamiento del lenguaje natural y la medicina. Su éxito radica en la habilidad de aprender directamente de los datos, ajustando millones de parámetros para realizar tareas de predicción, clasificación o regresión con una precisión notable.

Sin embargo, la creciente dependencia de estos modelos ha generado preocupaciones fundamentales en torno a la privacidad y la confidencialidad. El procesamiento de información sensible, como registros médicos, datos financieros o información personal, por parte de servicios de ML alojados en la nube o gestionados por terceros, expone a los usuarios a riesgos de fuga o uso indebido de sus datos. La necesidad de proteger la privacidad sin sacrificar la utilidad de los modelos de ML ha impulsado la investigación en técnicas de cómputo que preservan la privacidad.

En este contexto, el cifrado homomórfico emerge como una solución prometedora. El cifrado homomórfico permite realizar operaciones computacionales directamente sobre datos cifrados, de tal manera que el resultado de dichas operaciones, una vez descifrado, es idéntico al que se habría obtenido de realizar las mismas operaciones sobre los datos en texto plano. Esta propiedad única ofrece un mecanismo robusto para procesar información sensible en entornos no confiables, asegurando que los datos permanezcan cifrados durante todo el ciclo de procesamiento.

A pesar de su potencial, la aplicación práctica del cifrado homomórfico, especialmente en tareas computacionalmente intensivas como la inferencia de redes neuronales, ha estado limitada por su alto costo computacional. Los esquemas de cifrado homomórfico más avanzados, como CKKS (Cheon-Kim-Kim-Song), son capaces de manejar aritmética aproximada sobre números reales y complejos, lo cual los hace particularmente adecuados para las operaciones de punto flotante características del ML. No obstante, las implementaciones existentes de CKKS suelen depender de la CPU, lo que introduce cuellos de botella significativos cuando se trabaja con tensores grandes o modelos complejos.

La presente tesis aborda esta limitación proponiendo el desarrollo de una librería Python que aproveche la capacidad de procesamiento paralelo de las Unidades de Procesamiento Gráfico (GPU) para acelerar las operaciones tensoriales bajo el esquema de cifrado homomórfico CKKS. Nuestro objetivo es facilitar la inferencia privada de modelos de aprendizaje automático, ofreciendo una herramienta eficiente que combine la seguridad criptográfica con el rendimiento computacional necesario para aplicaciones del mundo real.

El resto de esta tesis se estructura en tres capítulos principales, cada uno diseñado para sentar las bases teóricas y justificar la propuesta:

\textbf{Capítulo 1: Fundamentos del Cifrado Homomórfico}
Este capítulo introduce los conceptos esenciales del cifrado homomórfico. Se exploran sus orígenes, la problemática del ruido inherente a estos esquemas, las diferencias entre cifrado parcialmente, algo y totalmente homomórfico, y las técnicas de gestión del ruido como el \textit{bootstrapping}. Se detallan los problemas matemáticos subyacentes, como aprendizjae con errores y aprendizaje con errores sobre anillos, y se presenta en profundidad el esquema CKKS, incluyendo sus mecanismos de codificación, generación de claves y operaciones homomórficas.

\textbf{Capítulo 2: Cifrado Homomórfico Aplicado al Machine Learning}
Aquí se examina la intersección entre el cifrado homomórfico y el aprendizaje automático. Se discute cómo las arquitecturas de redes neuronales, tanto profundas como convolucionales, pueden adaptarse para operar sobre datos cifrados. Se profundiza en la implementación de operaciones lineales clave (producto punto, multiplicación matriz-vector y convoluciones) en el dominio homomórfico, y se aborda el desafío de las funciones de activación no lineales y sus aproximaciones polinómicas. 

\textbf{Capítulo 3: Propuesta de Trabajo Especial de Grado}
Este capítulo presenta la propuesta central de esta tesis. Se explica el objetivo general y los objetivos específicos que guiarán el desarrollo de la librería. Se identifican las tecnologías clave a utilizar para la creacion de la libreria. Finalmente, se detalla un plan de actividades que describe las fases de implementación, pruebas y comparación de rendimiento con otras soluciones existentes.

